% =====================================================================
%  Stildefinitionen fuer die Faust-Blockdiagramme (Kapitel 5)
%  Darstellung angelehnt an die Ausgabe von "faust -svg":
%    - alle Primitive als Rechtecke (auch + - * ), nicht als Kreise
%    - Ein-/Ausgangsports auf den Blockkanten
%    - orthogonale Draehte mit Pfeilspitze in der Drahtmitte
%    - verschachtelte Rahmen fuer Kompositionsebenen
%    - rekursive Teilbloecke gespiegelt (Eingang rechts, Ausgang links)
%
%  EINBINDEN:  in die PRAEAMBEL der Arbeit:
%      % =====================================================================
%  Stildefinitionen fuer die Faust-Blockdiagramme (Kapitel 5)
%  Darstellung angelehnt an die Ausgabe von "faust -svg":
%    - alle Primitive als Rechtecke (auch + - * ), nicht als Kreise
%    - Ein-/Ausgangsports auf den Blockkanten
%    - orthogonale Draehte mit Pfeilspitze in der Drahtmitte
%    - verschachtelte Rahmen fuer Kompositionsebenen
%    - rekursive Teilbloecke gespiegelt (Eingang rechts, Ausgang links)
%
%  EINBINDEN:  in die PRAEAMBEL der Arbeit:
%      % =====================================================================
%  Stildefinitionen fuer die Faust-Blockdiagramme (Kapitel 5)
%  Darstellung angelehnt an die Ausgabe von "faust -svg":
%    - alle Primitive als Rechtecke (auch + - * ), nicht als Kreise
%    - Ein-/Ausgangsports auf den Blockkanten
%    - orthogonale Draehte mit Pfeilspitze in der Drahtmitte
%    - verschachtelte Rahmen fuer Kompositionsebenen
%    - rekursive Teilbloecke gespiegelt (Eingang rechts, Ausgang links)
%
%  EINBINDEN:  in die PRAEAMBEL der Arbeit:
%      % =====================================================================
%  Stildefinitionen fuer die Faust-Blockdiagramme (Kapitel 5)
%  Darstellung angelehnt an die Ausgabe von "faust -svg":
%    - alle Primitive als Rechtecke (auch + - * ), nicht als Kreise
%    - Ein-/Ausgangsports auf den Blockkanten
%    - orthogonale Draehte mit Pfeilspitze in der Drahtmitte
%    - verschachtelte Rahmen fuer Kompositionsebenen
%    - rekursive Teilbloecke gespiegelt (Eingang rechts, Ausgang links)
%
%  EINBINDEN:  in die PRAEAMBEL der Arbeit:
%      \input{abb/faust-diagramm-stile}
%
%  WICHTIG -- VOR DER ABGABE PRUEFEN:
%  Die innere Struktur von fi.tf2, fi.fir und an.fft ist aus den
%  Definitionen der Faust-Standardbibliothek (filters.lib, analyzers.lib)
%  rekonstruiert. Verifiziere sie gegen die von deinem Compiler tatsaechlich
%  erzeugten Diagramme:
%
%      faust -svg dsp/biquad.dsp   &&  xdg-open biquad-svg/process.svg
%      faust -svg dsp/fir32.dsp
%      faust -svg dsp/fft.dsp
%
%  Weicht etwas ab, gilt die Ausgabe von faust -svg -- sie stammt aus der
%  Version 2.83.1, die auch die gemessenen Klassen erzeugt hat.
% =====================================================================

% ---------------------------------------------------------------------
% (1) IN DIE PRAEAMBEL DER ARBEIT
% ---------------------------------------------------------------------
%
%   \usepackage{tikz}
%   \usetikzlibrary{positioning,calc,arrows.meta,fit,backgrounds,
%                   decorations.markings}
%   \usepackage{xcolor}
%   \usepackage{subcaption}      % nur fuer die FFT-Abbildung
%
% ---------------------------------------------------------------------
% (2) STILDEFINITIONEN -- ebenfalls in die Praeambel
% ---------------------------------------------------------------------

% Fuer reinen Schwarzweissdruck: die drei Fuellfarben auf "white" setzen.
\definecolor{fbxfill} {RGB}{214,230,246}   % Primitive und Bibliotheksbloecke
\definecolor{fuifill} {RGB}{252,235,203}   % Bedienelemente
\definecolor{fdelfill}{RGB}{247,221,221}   % Verzoegerungsglieder
\definecolor{fgrpfill}{RGB}{248,248,248}   % Kompositionsrahmen
\definecolor{fgrpline}{RGB}{130,130,130}

\tikzset{
  % Draht mit Pfeilspitze in der Drahtmitte (wie faust -svg)
  fw/.style       = {semithick,
                     decoration={markings, mark=at position #1 with
                       {\arrow{Latex[length=1.5mm,width=1.3mm]}}},
                     postaction={decorate}},
  fw/.default     = 0.6,
  fplain/.style   = {semithick},
  % Bloecke
  fb/.style       = {draw, semithick, fill=fbxfill, inner xsep=4pt,
                     inner ysep=2pt, minimum height=7mm, minimum width=9mm,
                     font=\small, align=center},
  fb2/.style      = {fb, minimum height=10mm},   % Block mit zwei Eingaengen
  fd/.style       = {fb, fill=fdelfill},        % Verzoegerungsglied
  fsum/.style     = {fb, minimum width=7mm, minimum height=7mm,
                     inner sep=1pt},            % Summierstelle
  fui/.style      = {draw, semithick, fill=fuifill, inner xsep=4pt,
                     inner ysep=2pt, minimum height=7mm, font=\footnotesize},
  % Kompositionsrahmen
  fgrp/.style     = {draw=fgrpline, thin, fill=fgrpfill, inner sep=4.5mm},
  fgrpin/.style   = {draw=fgrpline, thin, densely dashed, fill=white,
                     inner sep=2.4mm},
  fglab/.style    = {font=\tiny\itshape, text=fgrpline, inner sep=1pt},
  fsig/.style     = {font=\small, inner sep=1.5pt},
  flab/.style     = {font=\footnotesize, inner sep=2pt},
  fdot/.style     = {circle, fill, inner sep=0pt, minimum size=2.2pt},
}

% Kurzform fuer die beiden Eingangsports eines Blocks mit zwei Eingaengen
\newcommand{\pinA}[1]{($(#1.west)+(0,2.5mm)$)}   % oberer  Eingang
\newcommand{\pinB}[1]{($(#1.west)+(0,-2.5mm)$)}  % unterer Eingang

% Rahmenbeschriftung oberhalb der linken oberen Ecke
\newcommand{\fgname}[2]{\node[fglab, anchor=south west]
  at ([xshift=0.8mm,yshift=0.3mm]#1.north west) {#2};}

%
%  WICHTIG -- VOR DER ABGABE PRUEFEN:
%  Die innere Struktur von fi.tf2, fi.fir und an.fft ist aus den
%  Definitionen der Faust-Standardbibliothek (filters.lib, analyzers.lib)
%  rekonstruiert. Verifiziere sie gegen die von deinem Compiler tatsaechlich
%  erzeugten Diagramme:
%
%      faust -svg dsp/biquad.dsp   &&  xdg-open biquad-svg/process.svg
%      faust -svg dsp/fir32.dsp
%      faust -svg dsp/fft.dsp
%
%  Weicht etwas ab, gilt die Ausgabe von faust -svg -- sie stammt aus der
%  Version 2.83.1, die auch die gemessenen Klassen erzeugt hat.
% =====================================================================

% ---------------------------------------------------------------------
% (1) IN DIE PRAEAMBEL DER ARBEIT
% ---------------------------------------------------------------------
%
%   \usepackage{tikz}
%   \usetikzlibrary{positioning,calc,arrows.meta,fit,backgrounds,
%                   decorations.markings}
%   \usepackage{xcolor}
%   \usepackage{subcaption}      % nur fuer die FFT-Abbildung
%
% ---------------------------------------------------------------------
% (2) STILDEFINITIONEN -- ebenfalls in die Praeambel
% ---------------------------------------------------------------------

% Fuer reinen Schwarzweissdruck: die drei Fuellfarben auf "white" setzen.
\definecolor{fbxfill} {RGB}{214,230,246}   % Primitive und Bibliotheksbloecke
\definecolor{fuifill} {RGB}{252,235,203}   % Bedienelemente
\definecolor{fdelfill}{RGB}{247,221,221}   % Verzoegerungsglieder
\definecolor{fgrpfill}{RGB}{248,248,248}   % Kompositionsrahmen
\definecolor{fgrpline}{RGB}{130,130,130}

\tikzset{
  % Draht mit Pfeilspitze in der Drahtmitte (wie faust -svg)
  fw/.style       = {semithick,
                     decoration={markings, mark=at position #1 with
                       {\arrow{Latex[length=1.5mm,width=1.3mm]}}},
                     postaction={decorate}},
  fw/.default     = 0.6,
  fplain/.style   = {semithick},
  % Bloecke
  fb/.style       = {draw, semithick, fill=fbxfill, inner xsep=4pt,
                     inner ysep=2pt, minimum height=7mm, minimum width=9mm,
                     font=\small, align=center},
  fb2/.style      = {fb, minimum height=10mm},   % Block mit zwei Eingaengen
  fd/.style       = {fb, fill=fdelfill},        % Verzoegerungsglied
  fsum/.style     = {fb, minimum width=7mm, minimum height=7mm,
                     inner sep=1pt},            % Summierstelle
  fui/.style      = {draw, semithick, fill=fuifill, inner xsep=4pt,
                     inner ysep=2pt, minimum height=7mm, font=\footnotesize},
  % Kompositionsrahmen
  fgrp/.style     = {draw=fgrpline, thin, fill=fgrpfill, inner sep=4.5mm},
  fgrpin/.style   = {draw=fgrpline, thin, densely dashed, fill=white,
                     inner sep=2.4mm},
  fglab/.style    = {font=\tiny\itshape, text=fgrpline, inner sep=1pt},
  fsig/.style     = {font=\small, inner sep=1.5pt},
  flab/.style     = {font=\footnotesize, inner sep=2pt},
  fdot/.style     = {circle, fill, inner sep=0pt, minimum size=2.2pt},
}

% Kurzform fuer die beiden Eingangsports eines Blocks mit zwei Eingaengen
\newcommand{\pinA}[1]{($(#1.west)+(0,2.5mm)$)}   % oberer  Eingang
\newcommand{\pinB}[1]{($(#1.west)+(0,-2.5mm)$)}  % unterer Eingang

% Rahmenbeschriftung oberhalb der linken oberen Ecke
\newcommand{\fgname}[2]{\node[fglab, anchor=south west]
  at ([xshift=0.8mm,yshift=0.3mm]#1.north west) {#2};}

%
%  WICHTIG -- VOR DER ABGABE PRUEFEN:
%  Die innere Struktur von fi.tf2, fi.fir und an.fft ist aus den
%  Definitionen der Faust-Standardbibliothek (filters.lib, analyzers.lib)
%  rekonstruiert. Verifiziere sie gegen die von deinem Compiler tatsaechlich
%  erzeugten Diagramme:
%
%      faust -svg dsp/biquad.dsp   &&  xdg-open biquad-svg/process.svg
%      faust -svg dsp/fir32.dsp
%      faust -svg dsp/fft.dsp
%
%  Weicht etwas ab, gilt die Ausgabe von faust -svg -- sie stammt aus der
%  Version 2.83.1, die auch die gemessenen Klassen erzeugt hat.
% =====================================================================

% ---------------------------------------------------------------------
% (1) IN DIE PRAEAMBEL DER ARBEIT
% ---------------------------------------------------------------------
%
%   \usepackage{tikz}
%   \usetikzlibrary{positioning,calc,arrows.meta,fit,backgrounds,
%                   decorations.markings}
%   \usepackage{xcolor}
%   \usepackage{subcaption}      % nur fuer die FFT-Abbildung
%
% ---------------------------------------------------------------------
% (2) STILDEFINITIONEN -- ebenfalls in die Praeambel
% ---------------------------------------------------------------------

% Fuer reinen Schwarzweissdruck: die drei Fuellfarben auf "white" setzen.
\definecolor{fbxfill} {RGB}{214,230,246}   % Primitive und Bibliotheksbloecke
\definecolor{fuifill} {RGB}{252,235,203}   % Bedienelemente
\definecolor{fdelfill}{RGB}{247,221,221}   % Verzoegerungsglieder
\definecolor{fgrpfill}{RGB}{248,248,248}   % Kompositionsrahmen
\definecolor{fgrpline}{RGB}{130,130,130}

\tikzset{
  % Draht mit Pfeilspitze in der Drahtmitte (wie faust -svg)
  fw/.style       = {semithick,
                     decoration={markings, mark=at position #1 with
                       {\arrow{Latex[length=1.5mm,width=1.3mm]}}},
                     postaction={decorate}},
  fw/.default     = 0.6,
  fplain/.style   = {semithick},
  % Bloecke
  fb/.style       = {draw, semithick, fill=fbxfill, inner xsep=4pt,
                     inner ysep=2pt, minimum height=7mm, minimum width=9mm,
                     font=\small, align=center},
  fb2/.style      = {fb, minimum height=10mm},   % Block mit zwei Eingaengen
  fd/.style       = {fb, fill=fdelfill},        % Verzoegerungsglied
  fsum/.style     = {fb, minimum width=7mm, minimum height=7mm,
                     inner sep=1pt},            % Summierstelle
  fui/.style      = {draw, semithick, fill=fuifill, inner xsep=4pt,
                     inner ysep=2pt, minimum height=7mm, font=\footnotesize},
  % Kompositionsrahmen
  fgrp/.style     = {draw=fgrpline, thin, fill=fgrpfill, inner sep=4.5mm},
  fgrpin/.style   = {draw=fgrpline, thin, densely dashed, fill=white,
                     inner sep=2.4mm},
  fglab/.style    = {font=\tiny\itshape, text=fgrpline, inner sep=1pt},
  fsig/.style     = {font=\small, inner sep=1.5pt},
  flab/.style     = {font=\footnotesize, inner sep=2pt},
  fdot/.style     = {circle, fill, inner sep=0pt, minimum size=2.2pt},
}

% Kurzform fuer die beiden Eingangsports eines Blocks mit zwei Eingaengen
\newcommand{\pinA}[1]{($(#1.west)+(0,2.5mm)$)}   % oberer  Eingang
\newcommand{\pinB}[1]{($(#1.west)+(0,-2.5mm)$)}  % unterer Eingang

% Rahmenbeschriftung oberhalb der linken oberen Ecke
\newcommand{\fgname}[2]{\node[fglab, anchor=south west]
  at ([xshift=0.8mm,yshift=0.3mm]#1.north west) {#2};}

%
%  WICHTIG -- VOR DER ABGABE PRUEFEN:
%  Die innere Struktur von fi.tf2, fi.fir und an.fft ist aus den
%  Definitionen der Faust-Standardbibliothek (filters.lib, analyzers.lib)
%  rekonstruiert. Verifiziere sie gegen die von deinem Compiler tatsaechlich
%  erzeugten Diagramme:
%
%      faust -svg dsp/biquad.dsp   &&  xdg-open biquad-svg/process.svg
%      faust -svg dsp/fir32.dsp
%      faust -svg dsp/fft.dsp
%
%  Weicht etwas ab, gilt die Ausgabe von faust -svg -- sie stammt aus der
%  Version 2.83.1, die auch die gemessenen Klassen erzeugt hat.
% =====================================================================

% ---------------------------------------------------------------------
% (1) IN DIE PRAEAMBEL DER ARBEIT
% ---------------------------------------------------------------------
%
%   \usepackage{tikz}
%   \usetikzlibrary{positioning,calc,arrows.meta,fit,backgrounds,
%                   decorations.markings}
%   \usepackage{xcolor}
%   \usepackage{subcaption}      % nur fuer die FFT-Abbildung
%
% ---------------------------------------------------------------------
% (2) STILDEFINITIONEN -- ebenfalls in die Praeambel
% ---------------------------------------------------------------------

% Fuer reinen Schwarzweissdruck: die drei Fuellfarben auf "white" setzen.
\definecolor{fbxfill} {RGB}{214,230,246}   % Primitive und Bibliotheksbloecke
\definecolor{fuifill} {RGB}{252,235,203}   % Bedienelemente
\definecolor{fdelfill}{RGB}{247,221,221}   % Verzoegerungsglieder
\definecolor{fgrpfill}{RGB}{248,248,248}   % Kompositionsrahmen
\definecolor{fgrpline}{RGB}{130,130,130}

\tikzset{
  % Draht mit Pfeilspitze in der Drahtmitte (wie faust -svg)
  fw/.style       = {semithick,
                     decoration={markings, mark=at position #1 with
                       {\arrow{Latex[length=1.5mm,width=1.3mm]}}},
                     postaction={decorate}},
  fw/.default     = 0.6,
  fplain/.style   = {semithick},
  % Bloecke
  fb/.style       = {draw, semithick, fill=fbxfill, inner xsep=4pt,
                     inner ysep=2pt, minimum height=7mm, minimum width=9mm,
                     font=\small, align=center},
  fb2/.style      = {fb, minimum height=10mm},   % Block mit zwei Eingaengen
  fd/.style       = {fb, fill=fdelfill},        % Verzoegerungsglied
  fsum/.style     = {fb, minimum width=7mm, minimum height=7mm,
                     inner sep=1pt},            % Summierstelle
  fui/.style      = {draw, semithick, fill=fuifill, inner xsep=4pt,
                     inner ysep=2pt, minimum height=7mm, font=\footnotesize},
  % Kompositionsrahmen
  fgrp/.style     = {draw=fgrpline, thin, fill=fgrpfill, inner sep=4.5mm},
  fgrpin/.style   = {draw=fgrpline, thin, densely dashed, fill=white,
                     inner sep=2.4mm},
  fglab/.style    = {font=\tiny\itshape, text=fgrpline, inner sep=1pt},
  fsig/.style     = {font=\small, inner sep=1.5pt},
  flab/.style     = {font=\footnotesize, inner sep=2pt},
  fdot/.style     = {circle, fill, inner sep=0pt, minimum size=2.2pt},
}

% Kurzform fuer die beiden Eingangsports eines Blocks mit zwei Eingaengen
\newcommand{\pinA}[1]{($(#1.west)+(0,2.5mm)$)}   % oberer  Eingang
\newcommand{\pinB}[1]{($(#1.west)+(0,-2.5mm)$)}  % unterer Eingang

% Rahmenbeschriftung oberhalb der linken oberen Ecke
\newcommand{\fgname}[2]{\node[fglab, anchor=south west]
  at ([xshift=0.8mm,yshift=0.3mm]#1.north west) {#2};}
